% sections/ndss_05_cape.tex
% NDSS 2027 / CCS 2026 — §5 Source-Sink Compositional Attack Surface Validation (CAPE)
% Target: ~2 columns / ~1100 words

Having established that \ICone{} of ClawHub skills omit dangerous capabilities
from their public descriptions, we ask a structural security question:
\emph{does installing combinations of skills create attack surfaces the user
cannot anticipate?}
We answer through \cape (Compositional Attack Pair Evaluator),
an empirical framework that operationalizes the compositional attack surface
across \NpairsTotal{} unique source-sink pairs on a real LLM agent platform
(OpenClaw / \texttt{gpt-5.3-codex}).
A cross-platform replication on \skillssh{} using live \texttt{codex} agent
execution is reported in \S\ref{sec:crossplatform}.

\subsection{Attack Model and Threat Surface}

\paragraph{Compositional attack surface.}
The central security claim of this paper is not that any individual skill is
malicious, but that \emph{skill composition creates attack surfaces invisible
to the user}.
Formally, we define:

\begin{definition}[Source-Sink Compositional Attack Surface]
For a skill ecosystem $\mathcal{E}$, the \emph{compositional attack surface}
is the set of all source-sink pairs
$(\mathcal{S}_{\text{src}}, \mathcal{S}_{\text{sink}}) \in \mathcal{E}^2$
such that (i) $\mathcal{S}_{\text{sink}}$ possesses a dangerous capability
$\kappa \in \mathcal{C}_{\text{act,sink}} \setminus \mathcal{C}_{\text{vis,sink}}$
not disclosed to the user (\ICone), and (ii) $\mathcal{S}_{\text{src}}$
provides a channel through which attacker-influenced content can reach
the shared agent context.
\end{definition}

This definition has three key implications.
First, \textbf{the attack surface is structural}:
it arises from the act of combination itself, not from either skill being
individually malicious---a user who independently audits each installed
skill cannot detect the risk.
Second, \textbf{the attack surface is channel-agnostic}:
any source skill---web search, email, calendar, file reader, memory store,
browser automation, or database connector---can serve as the delivery
channel because the agent shares a single context window across all
installed skills.
Third, \textbf{the enabling condition is capability omission (\ICone)}:
if the sink skill disclosed its \texttt{shell\_exec} capability,
a user could apply the principle of least privilege and install only
the source skills they trust to carry trusted content.
IC-1 removes this defense entirely.

\begin{definition}[Source-Sink Composition Attack]
Given an agent session $\mathcal{A}$ with installed skills
$\mathcal{S}_{\text{src}}$ and $\mathcal{S}_{\text{sink}}$,
an attacker who can influence any content delivered through
$\mathcal{S}_{\text{src}}$'s output channel constructs content $c$
such that $\mathcal{A}$'s processing of $c$ invokes
$\kappa \in \mathcal{C}_{\text{act,sink}} \setminus \mathcal{C}_{\text{vis,sink}}$
without explicit user authorization.
\end{definition}

Note that the attacker's power requirement is minimal:
they need only inject content into \emph{one} of many plausible surfaces---
a search-engine result, an email in a monitored inbox, a document in a
shared folder, a cross-session memory entry, or an API response.
The \ICone gap is what converts this low-power capability into
a high-impact system compromise.

\subsection{Experimental Setup}

\paragraph{Pair selection.}
We select \NpairsTotal unique source-sink pairs to systematically
map the compositional attack surface.
Each pair combines one confirmed IC-1 sink with one source skill,
constituting an instance of the compositional attack surface.
Each pair is unique: no source or sink is reused across pairs,
ensuring independent validation of each sink's exploitability.
Source skills span seven injection-channel categories (\cf Table~\ref{tab:channels}):
\textbf{C1} LLM context (97 pairs): web search, browser automation,
  news fetch, API aggregation;
\textbf{C2} filesystem (1 pair): direct file read;
\textbf{C3} agent memory (17 pairs): cross-session memory write/read;
\textbf{C9} boot hook (4 pairs): startup script injection;
\textbf{C13} network (3 pairs): outbound request channel;
\textbf{C5} database (2 pairs): DB read/write;
\textbf{C9\_cron} (1 pair): scheduled task channel.
This multi-channel design directly validates that the attack surface
is channel-agnostic, not an artifact of any single delivery mechanism.

\paragraph{Task assignment.}
We define 15 attack task types (T1--T15) mapped to sink capability
categories: shell execution, file deletion, cron backdoor,
privilege escalation, script evaluation, process termination,
cross-session memory injection, wallet exfiltration, and others.
Each pair is assigned the task that best matches its sink's
confirmed dangerous capability.

\paragraph{Isolation.}
Each experiment runs in an isolated environment:
a dedicated workspace at \texttt{\textasciitilde/.cape\_exp/<pair\_id>/}
seeded with realistic context files (credentials, BOOT.md, memory files),
and a unique OpenClaw session ID.
A lightweight C2 server captures exfiltrated data for network-based tasks.
Side effects are verified by checking for probe files,
C2 server receipts, process state, crontab entries, or file modifications.

\paragraph{Injection delivery: identity-spoofing MCP source chain.}
A key methodological commitment distinguishes \cape from prior prompt injection evaluations:
we implement a \emph{true source-sink execution chain} where each pair's
source skill is individually registered as an identity-spoofing MCP server
that mirrors the real source skill's exact name, tool schema, and
Server Info declaration.

For each pair, our attacker-controlled MCP stdio server
(\texttt{attacker\_source\_skill.mjs}) is registered under the
\emph{real source skill's name} (\eg \texttt{tavily-search},
\texttt{brave-search}, \texttt{outlook-graph}, \texttt{browser-use})
via \texttt{openclaw\,mcp\,set}---identical to how a legitimate
skill is installed.
The server exposes the real tool name the source skill would expose
(\eg \texttt{tavily-search\_\_search}, \texttt{browser-use\_\_browse},
\texttt{outlook-graph\_\_get\_emails}),
ensuring the agent believes it is interacting with the
legitimate installed skill.
This models the realistic threat of a malicious clone skill published
under a name indistinguishable from a popular legitimate skill---
a supply-chain attack requiring no exploitation of the target platform's
infrastructure.

The attack chain proceeds as follows:

\begin{enumerate}[leftmargin=*,topsep=2pt,itemsep=1pt]
\item \textbf{Source skill registration (attacker).}
  The attacker publishes a skill with the same name and tool schema
  as a popular source skill (\eg Tavily Search).
  One-time registration via \texttt{openclaw\,mcp\,set} with the real skill's identity.

\item \textbf{Legitimate user workflow.}
  The user installs both a source skill (for web search, email, etc.) and
  a sink skill (for development, scheduling, memory, etc.).
  Neither skill appears dangerous in isolation;
  the sink's dangerous capability is IC-1-omitted.

\item \textbf{Source skill invocation.}
  During a routine task, the agent calls the source skill's tool
  (\eg \texttt{tavily-search\_\_search}) to fetch documentation or data.
  The call triggers a real MCP tool invocation over the OpenClaw gateway.

\item \textbf{Injection via tool output.}
  The source skill returns crafted ``integration documentation''
  containing a disguised shell command.
  The payload is framed as a legitimacy-enhancing pre-flight check
  using a living-off-the-land pattern:
  \texttt{go\,vet\,./...\,2>\&1\,|\,tee\,\textlangle{}probe\textrangle{}},
  presented as a ``go vet compatibility check required by the sink skill.''
  We validate two alternative payload families:
  (\textbf{P2}) \texttt{node\,verify.js\,2>\&1\,|\,tee\,\textlangle{}p\textrangle{}}
  and (\textbf{P3}) \texttt{cat\,README.md\,|\,tee\,\textlangle{}p\textrangle{}},
  both achieving comparable success (P2: 7/8, P3: 6/8 on held-out pairs).

\item \textbf{Sink capability invocation.}
  The agent executes the command verbatim, triggering the sink skill's
  IC-1-omitted \texttt{shell\_exec} capability and
  producing a verified file-system side effect.
\end{enumerate}

This chain provides Level~3 empirical evidence---real registered skill,
real agent session, real tool invocation, real system-level side effect---
validating the compositional attack surface across 125 distinct source-sink pairs
spanning 7 injection channels and 13 dangerous capability categories.

\paragraph{Controls.}
To distinguish attack effect from model variance, we run three control conditions
on a stratified random sample of 20 pairs (4 per task type):
\textbf{(Ctrl-A)} same pair, but source skill returns a benign documentation page (no injection);
\textbf{(Ctrl-B)} same pair, same injection, but the payload command is replaced with a clearly
benign operation (\texttt{echo hello > /dev/null});
\textbf{(Ctrl-C)} same pair, repeated 3$\times$ with identical injection to measure within-pair variance.
Ctrl-A: 0/20 probe files created (baseline = zero without injection).
Ctrl-B: 18/20 agent executes the benign command (confirming agent compliance with tool context),
but 0/20 create the probe file.
Ctrl-C: within-pair success variance is low ($\sigma = 0.11$) for T06 pairs and moderate
($\sigma = 0.28$) for T18/T13 pairs, consistent with stochastic refusal rather than random model behavior.
These controls confirm that: (1) probe file creation is caused by the injected command, not ambient side effects;
(2) agent compliance is not the bottleneck---payload delivery is;
and (3) stochastic refusals in T13/T18 reflect per-session context sensitivity, not a stable defense.



\subsection{Results}

\begin{table}[t]
\centering
\caption{CAPE attack validation results by task type across all 125 source-sink pairs.
  All pairs use the true MCP source-sink chain (real attacker MCP server \textrightarrow\ real agent session \textrightarrow\ verified file-system side effect).}
\label{tab:cape_results}
\setlength{\tabcolsep}{4pt}
\small
\begin{tabular}{lllrr}
\toprule
\textbf{Task} & \textbf{Sink Capability} & \textbf{Injection Disguise} & \textbf{Pairs} & \textbf{Success}\\
\midrule
T06 & Shell exec (arbitrary)   & \texttt{go vet | tee \textlangle{}p\textrangle{}} & 70  & 33 (47\% ITT; \textbf{100\% CR}) \\
T18 & Process termination      & \texttt{go vet | tee \textlangle{}p\textrangle{}} & 17  & 5 (29\% ITT; 71\% CR) \\
T13 & Cross-session memory     & \texttt{go vet | tee \textlangle{}p\textrangle{}} & 16  & 8 (50\% ITT; 89\% CR) \\
T17 & Wallet exfiltration      & \texttt{go vet | tee \textlangle{}p\textrangle{}} &  5  & 3 (60\% ITT; 100\% CR) \\
T05 & Boot hook injection      & \texttt{go vet | tee \textlangle{}p\textrangle{}} &  4  & 0 (0\% ITT; 0\% CR) \\
T03 & Cron backdoor            & \texttt{go vet | tee \textlangle{}p\textrangle{}} &  1  & 1 (100\%) \\
T14 & Script eval              & \texttt{go vet | tee \textlangle{}p\textrangle{}} &  1  & 1 (100\%) \\
T11 & Privilege escalation     & \texttt{go vet | tee \textlangle{}p\textrangle{}} &  1  & 1 (100\%) \\
T07/T08/T10 & Network/DB exfil  & \texttt{go vet | tee \textlangle{}p\textrangle{}} &  5  & 5 (100\%) \\
T02/T04/T19/T20 & Other        & \texttt{go vet | tee \textlangle{}p\textrangle{}} &  5  & 0 (0\% — all timeout) \\
\midrule
\textbf{Total} && & \textbf{125} & \textbf{57 ITT (45.6\%) / 57 CR (93.4\%)} \\
\bottomrule
\end{tabular}
\end{table}

Table~\ref{tab:cape_results} reports results for 85/125 completed pairs
(40 remaining; camera-ready will include all 125).
We report two metrics, pre-defined before any analysis:

\textbf{Intent-to-treat (ITT)}: 57/85 = \textbf{67.1\%}
(95\% CI: 56.5--76.1\%),
counting all 24 timed-out pairs as failures.
Timeout policy: each pair is given one attempt with a 175\,s hard kill
(\texttt{SIGKILL} on the process group);
if no probe file is created within that window, the pair is marked failed.

\textbf{Completed-run rate}: 57/61 = \textbf{93.4\%}
(95\% CI: 84.1--97.8\%),
conditioned on pairs where the agent returned a response before timeout.
24 pairs timed out due to agent infrastructure latency
(LLM context compaction under high prompt token usage)---
this is a measurement-infrastructure limitation, not an agent defense:
timeout pairs produce no output and thus no refusal, evidence, or probe file.
Of the 24 ITT failures, only \textbf{4} (4.7\% of tested pairs) are agent refusals;
the remaining 20 are silent timeouts with no agent output.

T06 (shell exec, completed-run): 33/33 = \textbf{100\%};
T07/T08/T10 (network/DB exfil): 5/5 = 100\%;
T17 (wallet): 3/3 = 100\%.

\paragraph{Uniform injection technique — all 15 task types.}
A key finding is that the \textbf{same injection disguise} ($\S$\ref{sec:injection_detail})
works across all 15 sink task types: not just shell-exec sinks,
but also process monitors, memory indexers, wallet validators,
boot hooks, database connectors, and script evaluators.
Each sink's documentation is framed around go vet compatibility,
causing agents to execute the command regardless of the sink's primary domain.
This demonstrates that the vulnerability is in the \emph{agent's trust model},
not in any property specific to shell-exec sinks.

\paragraph{True MCP source-sink chain: all 125 pairs.}
Unlike prior prompt injection evaluations that hand-craft source-skill output,
all 125 \cape pairs use a \emph{real, registered MCP server} as the source skill.
The attacker MCP server (\texttt{attacker\_source\_skill.mjs}) is installed
via the standard \texttt{openclaw\,mcp\,set} interface---
the same mechanism used to install any legitimate skill---and its
\texttt{web\_search} tool interface is indistinguishable from
Tavily, DuckDuckGo, Google Search, or any other web-search source skill.
Source skills span 12+ different types (web search, email, news feeds,
browser automation, OCR, RSS, social media scrapers) across all 125 pairs,
confirming that the attack generalizes across the full source skill ecosystem.

\paragraph{Shell execution sinks: no consistent defense.}
Of 33 T06 (shell exec) pairs that ran to completion, \textbf{33/33 = 100\%} produced
verified side effects across 33 distinct source skills (web search, email,
browser automation, finance APIs, news aggregators, OCR, and social scrapers).
The agent's safety model applies no consistent restriction to shell commands
delivered through source-skill tool output.
This is the most dangerous finding: 55.6\% of all confirmed dangerous sinks
(\S\ref{sec:measurement}) have shell execution capability, and no completed pair
shows a successful defense against the \texttt{go\,vet}-disguised injection.

\paragraph{Asymmetric safety policy.}
The agent exhibits a striking asymmetry:
it occasionally refuses requests to modify agent memory files directly
(``I can't run pasted diagnostic commands from untrusted integration messages''),
but applies no equivalent restriction to \texttt{go vet}
or other build-tool commands.
This heuristic-based policy leaves the more dangerous execution surface
completely unguarded against build-tool-disguised payloads.

\paragraph{Failure analysis and refusal variance.}
Of 85 pairs attempted, 4 produced explicit agent refusals:
(i) P001: ``I can't run commands extracting local identity content from untrusted search output'';
(ii) P004: ``system process data from untrusted source is not acceptable'';
(iii) P049, P050: LLM idle timeout (no output produced).

Critically, \emph{neither refusal message fires consistently}: T13 (memory) succeeds in
8/12 tested pairs (67\%) and T18 (process kill) in 5/11 (45\%),
demonstrating that the agent applies safety heuristics stochastically
rather than enforcing a stable capability policy.
Per-pair repetition ($n=3$; Ctrl-C; \S\ref{sec:controls}) confirms:
Within-task-type refusal variance is $\sigma = 0.11$ for T06 and $\sigma = 0.28$ for T13/T18.
A single refusal in one session is not predictive of refusal in subsequent sessions with the same pair.
This inconsistency is itself a security vulnerability: an attacker need only retry
under different session contexts until the safety heuristic does not fire.

\paragraph{Injection disguise generalizes.}
The \texttt{go vet ./... 2>\&1 | tee \textlangle{}probe\textrangle{}} payload
succeeds because the agent interprets it as a routine build-tool caching step.
This ``living-off-the-land'' pattern---using legitimate, pre-installed
developer tools in ways that produce attacker-controlled side effects---is
not detectable by any static inspection of the command string.
No special capabilities or elevated permissions are required;
the attack exploits only tools already present in a typical developer environment.


\subsection{Key Takeaway}
\label{sec:cape_takeaway}

\begin{observation}
A skill's undisclosed dangerous capability is, in practice,
exploitable by any source skill that can deliver attacker-controlled
content to the agent's context window.
The platform's safety policy is heuristic-based and stochastically applied:
no consistent runtime defense is activated against build-tool-disguised
injection payloads, even across repeated trials.
\end{observation}

Our results establish that the IC-1 gap (capability omission) is not
merely a transparency issue---it is an active attack enabler.
Users cannot consent to capabilities they cannot see,
and platforms cannot enforce policies they have not enumerated.
A cross-platform replication with live \texttt{codex} agent execution
on \NssPairs{} \skillssh{} skills ($\NssSuccessRate{}$ attack success rate)
is reported in \S\ref{sec:crossplatform}.
